\section{Fundamentos analíticos y geométricos de la localización}
\label{sec:ruta-biblioteca-local-drive}

La localización de un atractor combina operaciones que pertenecen a teorías
distintas. El álgebra lineal determina equilibrios y modos locales; Fourier
proporciona una aproximación armónica; la teoría de ecuaciones diferenciales
permite interpretar la trayectoria; y la geometría de cuencas relaciona esa
trayectoria con otras respuestas del sistema. En orden fraccionario se añade
la memoria: el valor actual del estado no contiene, en general, toda la
información necesaria para continuar una solución.

\subsection{Dependencias entre los fundamentos}
\label{sec:ruta-unica-estudio}

El resolvente de una matriz exige conocer su espectro; una transferencia
fraccionaria añade la elección de una rama compleja; y el balance armónico
usa después ambas construcciones. Un mapa de Poincaré presupone existencia
y unicidad del flujo, además de transversalidad. Estas dependencias explican
por qué una identidad algebraica y una conclusión sobre la cuenca requieren
argumentos diferentes.

\begin{longtable}{@{}L{0.21\linewidth}L{0.34\linewidth}L{0.37\linewidth}@{}}
\caption{Fundamentos y problemas de localización.}
\label{tab:ruta-unica-biblioteca}\\
\toprule
Fundamento & Construcción & Aplicación\\
\midrule
\endfirsthead
\toprule
Fundamento & Construcción & Aplicación\\
\midrule
\endhead
Álgebra lineal & Jacobiano, espectro y resolvente. & Equilibrios, estabilidad y transferencia.\\
Análisis de operadores & Integrales beta y transformadas. & Volterra y respuesta de Caputo.\\
Fourier y variable compleja & Rama de $s^q$ y convolución armónica. & Semilla y balance con sesgo.\\
Teoría de EDO & Existencia y ecuación variacional. & Flujo y sensibilidad inicial.\\
Geometría dinámica & Retorno y variedades invariantes. & Ciclos y separatrices.\\
Topología & Grado y conjuntos aislantes. & Existencia de dinámica invariante.\\
Memoria & Perfiles de historia. & Continuación causal y cuencas funcionales.\\
Análisis numérico & Refinamiento y comparación independiente. & Persistencia de los resultados.\\
\bottomrule
\end{longtable}

El cálculo multivariable y las EDO se desarrollan en
Kantorovich~\cite{Kantorovich2022Fundamentals}, pp.\ 243--760; el álgebra
espectral, el análisis complejo, Fourier y Laplace se reúnen en
\cite{Kantorovich2024Advanced}, pp.\ 1--717. Los ejercicios de
Cicogna~\cite{Cicogna2020Exercises}, pp.\ 57--127, relacionan transformadas,
distribuciones y representación frecuencial. Para la dinámica no lineal,
Wimberger~\cite{Wimberger2022Nonlinear}, pp.\ 9--152, conecta sistemas
disipativos y formulación geométrica.

\subsection{Definición, inicialización y respuesta fraccionaria}

\begin{table}[H]
\centering
\caption{Perspectivas analíticas del cálculo fraccionario.}
\label{tab:ruta-frac-books}
\small
\begin{tabularx}{\textwidth}{@{}L{0.24\textwidth}L{0.29\textwidth}Y@{}}
\toprule
Fuente & Desarrollo analítico & Relación con el problema dinámico \\
\midrule
Oliveira \cite{Oliveira2019SolvedFractional} &
Funciones especiales (pp.\ 17--67), Mittag--Leffler (69--113), transformadas
integrales (115--167), derivadas fraccionarias (169--222) y aplicaciones
(223--307). &
Es la fuente de trabajo algebraico: Gamma, Mittag--Leffler, Laplace,
funciones de Green, ecuaciones integrales de Volterra y ejercicios con
solución. \\
Herrmann \cite{Herrmann2014Fractional} &
Operadores fraccionarios y semigrupo en el orden (caps.\ 3 y 5), oscilador
fraccionario (cap.\ 6), no localidad y memoria (cap.\ 8), y formalismos de
simetría, tensores y campos (caps.\ 13, 16 y 23--25). &
Ofrece intuición física y puentes geométricos. Sus ``grupos fraccionarios'' y
``fractors'' son formalismos específicos; no se presentan como geometría
universal de sistemas Caputo. \\
Das \cite{Das2020Kindergarten} &
Integración (pp.\ 47--99), derivadas (100--154), inicialización y Laplace
generalizada (210--254), ecuaciones diferenciales e integrales fraccionarias
(255--343) y ramas complejas en el apéndice E (466--470). &
Es un puente conceptual amplio, útil para terminal inferior, inicialización,
sistemas vectoriales y ramas de \(s^q\). Su tono deliberadamente informal
requiere respaldar los teoremas con una fuente analítica rigurosa. \\
\bottomrule
\end{tabularx}
\end{table}

Los números anteriores son páginas impresas. La correspondencia con la
paginación digital de los ejemplares consultados se presenta en la
\cref{tab:ruta-frac-pdf-pages}.

\begin{longtable}{@{}L{0.21\textwidth}L{0.35\textwidth}L{0.35\textwidth}@{}}
\caption{Correspondencia entre tema, página impresa y página PDF.}
\label{tab:ruta-frac-pdf-pages}\\
\toprule
Fuente & Tema y página impresa & Página del PDF \\
\midrule
\endfirsthead
\toprule
Fuente & Tema y página impresa & Página del PDF \\
\midrule
\endhead
Oliveira & Mittag--Leffler, pp.\ 69--113 & PDF 81--125 \\
Oliveira & Transformadas, pp.\ 115--167 & PDF 126--178 \\
Oliveira & Derivadas fraccionarias, pp.\ 169--222 & PDF 179--232 \\
Oliveira & Volterra/Green: ejercicios 33, 35, 36 y 38,
pp.\ 277, 280, 281 y 283 & PDF 287, 290, 291 y 293 \\
Das & Riemann--Liouville y Caputo, \S\S3.7--3.10,
pp.\ 119--125 & PDF 156--162 \\
Das & Inicialización, cap.\ 5, pp.\ 210--254 & PDF 247--291 \\
Das & Sistemas vectoriales, \S7.13, p.\ 322 y \S7.17, p.\ 335 &
PDF 359 y 372 \\
Das & Ramas complejas y superficies de Riemann, apéndice E,
pp.\ 466--470 & PDF 503--507 \\
Herrmann & Definiciones y operadores, \S3.4.2 y \S\S5.2--5.3,
pp.\ 26 y 45--52 & PDF 47 y 66--73 \\
Herrmann & Interpretación geométrica elemental, \S5.7, p.\ 62 & PDF 83 \\
Herrmann & Semigrupo en el orden, \S5.9.1, p.\ 66 & PDF 87 \\
Herrmann & No localidad y memoria, cap.\ 8, pp.\ 87--106 & PDF 108--127 \\
Herrmann & Tensores, campos y gauge, caps.\ 23--25,
pp.\ 353--378 & PDF 374--399 \\
\bottomrule
\end{longtable}

\subsection{Composición en el orden y evolución temporal}

Para $\alpha,\beta>0$ definimos
\[
 I_{t_0}^{\alpha}f(t)=\frac1{\Gamma(\alpha)}
 \int_{t_0}^t(t-u)^{\alpha-1}f(u)\,du.
\]
Si $f$ es continua sobre el intervalo finito considerado, el integrando
iterado es absolutamente integrable. Fubini permite intercambiar el orden
en el triángulo $t_0\le v\le u\le t$:
\[
 I_{t_0}^\alpha I_{t_0}^\beta f(t)=
 \frac1{\Gamma(\alpha)\Gamma(\beta)}
 \int_{t_0}^t f(v)\int_v^t
 (t-u)^{\alpha-1}(u-v)^{\beta-1}\,du\,dv.
\]
Con $u=v+(t-v)r$, $du=(t-v)dr$, la integral interior es
\[
 (t-v)^{\alpha+\beta-1}B(\beta,\alpha)
 =(t-v)^{\alpha+\beta-1}
 \frac{\Gamma(\alpha)\Gamma(\beta)}{\Gamma(\alpha+\beta)}.
\]
Por tanto,
\begin{equation}
 I_{t_0}^\alpha I_{t_0}^\beta=I_{t_0}^{\alpha+\beta}.
 \label{eq:library-order-semigroup}
\end{equation}
Esta identidad compone órdenes de integración, no tiempos de evolución.
La diferencia se ve en ${}^{\mathrm C}D_0^q x=\lambda x$, cuya solución
es $x(t)=E_q(\lambda t^q)x_0$. Para $0<q<1$, $\lambda\ne0$ y $h$ pequeño,
\begin{align*}
 E_q(\lambda(2h)^q)&=1+\frac{2^q\lambda h^q}{\Gamma(q+1)}+O(h^{2q}),\\
 E_q(\lambda h^q)^2&=1+\frac{2\lambda h^q}{\Gamma(q+1)}+O(h^{2q}).
\end{align*}
Como $2^q\ne2$, ambas expresiones difieren ya en el primer término no
constante. En consecuencia, la identidad
\begin{equation}
 \Phi(t+s,x_0)=\Phi\bigl(t,\Phi(s,x_0)\bigr)
 \label{eq:library-false-time-semigroup}
\end{equation}
no describe, en general, la evolución proyectada sobre el estado puntual de
Caputo. Reiniciar desde $x(s)$ omite la contribución de la historia anterior
a $s$, desarrollada en \cref{sec:geometria-topologia-caputo}.

La regla de cadena ordinaria presenta otra diferencia. Con $x(t)=t$ y
$g(x)=x^2$,
\[
 {}^{\mathrm C}D_0^q[t^2]=\frac{2t^{2-q}}{\Gamma(3-q)},\qquad
 g'(t){}^{\mathrm C}D_0^q[t]=\frac{2t^{2-q}}{\Gamma(2-q)}.
\]
Puesto que $\Gamma(3-q)=(2-q)\Gamma(2-q)$, las fórmulas sólo coinciden
en el límite entero $q=1$. Los cambios de coordenadas y las semillas FPP--A
se construyen por ello desde Volterra, sin suponer esa regla de cadena.

\subsection{Ramas complejas y transferencia}

Definimos $s^q=\exp(q\operatorname{Log}s)$ sobre
$\mathbb C\setminus(-\infty,0]$, con $\arg s\in(-\pi,\pi)$. Si
$s=\ii\omega$ y $\omega>0$,
\begin{equation}
 (\ii\omega)^q=\omega^q\exp\!\left(\ii q\frac\pi2\right).
 \label{eq:library-principal-branch}
\end{equation}
Sus partes real e imaginaria son $\omega^q\cos(q\pi/2)$ y
$\omega^q\sin(q\pi/2)$; para $q=1$ recuperamos $\ii\omega$. A ambos
lados del corte negativo, los límites son $r^qe^{\pm\ii q\pi}$,
$r>0$. La elección de rama se mantiene fija durante la continuación.
La superficie de Riemann de esta potencia y la topología de una cuenca
son objetos diferentes.

La transferencia $C(s^qI-A)^{-1}B$ existe donde el resolvente está definido.
Su signo depende del convenio de realimentación, de modo que el cierre debe
conservar simultáneamente la rama y la definición de $W_+$ o $W_-$.
Las derivaciones concretas se presentan en
\cref{sec:lure_transferencia,sec:machado-fdf}.

\subsection{Geometría de EDO y familias con atractores ocultos}

Hirsch--Smale--Devaney
\cite{HirschSmaleDevaney2013} conecta las construcciones utilizadas en los
modelos estudiados:
\begin{enumerate}
 \item capítulos 1--6: ecuaciones y sistemas lineales, retratos y
 clasificación;
 \item capítulos 7--9: existencia, dependencia continua, ecuación variacional,
 equilibrios y técnicas globales;
 \item capítulo 10: secciones locales, mapa de Poincaré, conjuntos límite y
 Poincaré--Bendixson;
 \item capítulo 12: circuitos, Liénard, Van der Pol y bifurcación de Hopf;
 \item capítulo 14: sistema y atractor de Lorenz;
 \item capítulo 15: dinámica discreta, caos, dinámica simbólica, corrimiento y
 conjunto de Cantor;
 \item capítulo 16: Shilnikov, herradura, doble scroll, homoclinicidad y una
 exploración del circuito de Chua;
 \item capítulo 17: prueba de existencia, unicidad, dependencia continua y
 diferenciabilidad del flujo.
\end{enumerate}

Hoppensteadt \cite{Hoppensteadt2000Chaotic} añade una ruta aplicada que encaja
con dos sistemas concretos: ecuaciones angulares y cilindro del PLL en
\S2.2; sistemas disipativos y PLL en \S2.4; mapas estroboscópicos en \S2.5;
variedades estable/inestable en \S3.2.3; estabilidad orbital y retorno de
Poincaré en \S3.5; y bifurcación, función implícita y métodos topológicos en el
capítulo 4. La identificación de fases módulo $2\pi$ forma parte del estado del PLL;
una distancia entre fases sin envolver puede separar artificialmente dos
representantes del mismo punto del cilindro.

El volumen de Pham et al. \cite{PhamEtAl2018HiddenAttractors} cumple otra
función. Su índice separa una parte de sistemas con atractores autoexcitados y
otra de sistemas con atractores ocultos. La segunda incluye órbitas
periódicas, toros invariantes y caos en sistemas cuadráticos sin equilibrio,
un oscilador Duffing memristivo, Rikitake, sistemas de cuatro y cinco
dimensiones y sincronización. Es un catálogo de modelos y antecedentes, no un
teorema general de ocultedad. La comparación entre esos modelos requiere fijar sus ecuaciones, sus
parámetros y sus condiciones iniciales; la estructura de los equilibrios
orienta después el estudio de la cuenca.

\begin{itemize}
  \item \emph{Chaotic Systems with Multistability and Hidden Attractors}
        \cite{WangKuznetsovChen2021Multistability}: sistemas con equilibrios
        estables, sin equilibrios, con curvas o superficies de equilibrios
        (pp.\ 29--146); sistemas hipercaóticos, fraccionarios, memristivos y
        jerk con atractores ocultos (pp.\ 149--308); y multistabilidad
        simétrica o asimétrica (pp.\ 311--375).
  \item \emph{Systems with Hidden Attractors}
        \cite{PhamVolosKapitaniak2017Systems}: clasificación compacta por
        equilibrios estables (pp.\ 21--35), infinitos equilibrios (37--50),
        ausencia de equilibrios (51--63), sincronización (65--77) y
        realización circuital (79--102).
  \item \emph{Nonlinear Dynamics} \cite{DatserisParlitz2022Nonlinear}:
        dinámica continua, definición y medición de caos, bifurcaciones,
        entropía, dimensión fractal, coordenadas de retardo y análisis de
        series de tiempo (pp.\ 1--119). Su código Julia es un complemento
        natural de la reproducción con \juliads, pero las hipótesis de cada algoritmo determinan qué diagnóstico puede
        inferirse de sus resultados.
  \item \emph{Elements of Applied Bifurcation Theory}
        \cite{Kuznetsov2023Bifurcation}: equivalencia topológica y estabilidad
        estructural (pp.\ 47--87), bifurcaciones de equilibrios y ciclos
        (89--228), órbitas homoclínicas y heteroclínicas (229--279) y análisis
        numérico de bifurcaciones (569--655).
\end{itemize}

\subsection{Sensibilidad, recurrencia y cuencas}
\label{sec:lecturas-por-procedimiento}

\begin{longtable}{@{}L{0.24\linewidth}L{0.34\linewidth}L{0.34\linewidth}@{}}
\caption{Información complementaria en el estudio de una respuesta.}
\label{tab:lectura-procedimientos}
\label{tab:drive-complementarios}
\label{tab:drive-isbn-cerrado}\\
\toprule
Objeto & Construcción & Interpretación\\
\midrule
\endfirsthead
\toprule
Objeto & Construcción & Interpretación\\
\midrule
\endhead
Equilibrios & Reducción escalar, pertenencia regional y espectro. & Estabilidad local y acceso desde su vecindad.\\
Balance armónico & Sesgo, proyección de Fourier y residual físico. & Condición inicial aproximada.\\
Continuación & Homotopía y ecuación integral con historia única. & Transporte de la respuesta.\\
Puntos perpetuos & $JF=0$, $F\ne0$ y desarrollo de arranque. & Cancelación local y destino posterior.\\
Exponentes finitos & Separación de trayectorias en una ventana. & Sensibilidad y predictibilidad local.\\
Secciones & Retornos y dispersión de cruces. & Ciclo o nube no periódica.\\
Variedades invariantes & Separatrices y condiciones de acceso. & Coexistencia y fronteras de cuenca.\\
Topología & Grado y conjuntos aislantes. & Existencia bajo hipótesis de continuidad y aislamiento.\\
\bottomrule
\end{longtable}

Una tasa de separación positiva durante una ventana puede pertenecer a un
transitorio. Una figura irregular tampoco determina por sí sola invariancia,
transitividad o entropía positiva. La combinación de trayectoria, retornos,
variaciones y cuenca relaciona cada diagnóstico con el objeto que mide.

En las EDO, existencia y unicidad permiten trasladar una condición inicial
a lo largo de su órbita. En Caputo, trasladar sólo su valor puntual define
en general otra solución. Dos historias con el mismo valor actual pueden
conservar contribuciones distintas del pasado al futuro. Un retorno
proyectado sobre $x$ puede perder la univocidad que tendría sobre un espacio
de perfiles adecuadamente definido.

Estas diferencias organizan la metodología: las identidades de realización
y balance proporcionan semillas; los integradores aproximan el PVI declarado;
y los diagnósticos describen sus soluciones en los intervalos estudiados.
La geometría relaciona después esas soluciones con invariancia, atracción
y fronteras de cuenca.
